\begin{frame}{손으로 한 번: ``수식적 정당화''의 좋은/나쁜 예}
  {\footnotesize 설정: 어떤 팀이 CliffWalking에서 ``Q-learning이 SARSA보다
  최종 평가 점수가 더 높았다 ($-13$ vs $-17$)''라는 결과를 얻었다.}
  \vskip0.5em
  \begin{alertblock}{나쁜 설명}
    ``Q-learning이 SARSA보다 점수가 더 높았다. Q-learning이 더 좋은
    알고리즘인 것 같다.''
  \end{alertblock}
  \vskip0.3em
  \begin{block}{좋은 설명 (요약)}
    학습 도중(탐험 포함)은 SARSA가 더 높았다 ($-31$ vs $-53$) --- Q-learning은
    \textbf{목표 정책(항상 최선)의 가치를 배우는 off-policy} 방법이라 절벽
    옆 최단 경로를 학습했고, 그 경로는 탐험 실수 한 번이 $-100$으로
    이어진다. SARSA는 \textbf{실제 행동 방식(탐험 포함)까지 감안하는
    on-policy}라 안전한 경로를 선호. 탐험 없이 평가하면 Q-learning 경로가
    더 짧아 $-13$ vs $-17$이지만, 이건 ``더 좋은 알고리즘''이 아니라
    \textbf{``두 알고리즘이 애초에 다른 질문에 답하고 있다''}.
  \end{block}
  \vskip0.4em
  \kq{같은 숫자를 말해도, 좋은 설명은 숫자 뒤의 \textbf{알고리즘의
  설계 차이}를 짚어낸다 --- 이것이 ``수식적 정당화''의 실제 의미.}
\end{frame}
